\section{Sliding Window Median}
\label{sec:cses-sw-median}

\problemstatement{Given an array of $n$ integers and a window size $k$,
output the median of every contiguous window (for even $k$, the lower of
the two middle values).}

\begin{patternbox}
The median needs the middle element under both insertion and deletion. Two
balanced \textbf{multisets} -- a \emph{low} half and a \emph{high} half --
split at the median do exactly this: keep the low half's size equal to (or
one more than) the high half, and the median is the low half's maximum.
\end{patternbox}

\subsection*{Two balanced halves}

Maintain \textit{low} (a multiset of the smaller elements) and \textit{high}
(the larger elements) with the invariant $|\textit{low}|=\lceil m/2\rceil$
for a window of size $m$. Insertion routes the new value to the correct half
and then rebalances by moving one element across if sizes drift. Deletion
(when the left element leaves) erases it from whichever half holds it -- a
multiset supports erase-by-value in $O(\log k)$ -- and rebalances. The
median is always $\ast\textit{low}.\text{rbegin}()$ (its largest element).

\begin{ideabox}[Balanced Multisets Give Erasable Heaps]
Ordinary heaps cannot delete an arbitrary element, which sliding windows
require. Two multisets behave like two heaps that \emph{also} support
$O(\log k)$ removal of a specific value, keeping the median at the boundary
between the halves.
\end{ideabox}

\subsection*{Algorithm}

\begin{enumerate}
  \item Insert into \textit{low} or \textit{high} by comparing with the
        current median; rebalance so $|\textit{low}|-|\textit{high}|
        \in\{0,1\}$.
  \item On slide, erase the departing value from its half and rebalance.
  \item Output $\ast\textit{low}.\text{rbegin}()$ per window.
\end{enumerate}

\subsection*{Dry run}

\dryrun{$a=[2,4,1,3]$, $k=3$.}

\begin{itemize}
  \item $[2,4,1]$ sorted $1,2,4$: low $=\{1,2\}$, high $=\{4\}$; median
        $2$.
  \item Slide to $[4,1,3]$ sorted $1,3,4$: erase $2$, insert $3$; low
        $=\{1,3\}$, high $=\{4\}$; median $3$.
  \item Output $2,3$.
\end{itemize}

\subsection*{C++ implementation}

\begin{lstlisting}
#include <bits/stdc++.h>
using namespace std;

int main() {
    int n, k;
    cin >> n >> k;
    vector<int> a(n);
    for (int& x : a) cin >> x;

    multiset<int> low, high;                  // low: smaller half
    auto rebalance = [&]() {
        while (low.size() > high.size() + 1) {
            high.insert(*low.rbegin());
            low.erase(prev(low.end()));
        }
        while (low.size() < high.size()) {
            low.insert(*high.begin());
            high.erase(high.begin());
        }
    };
    auto insert = [&](int x) {
        if (low.empty() || x <= *low.rbegin()) low.insert(x);
        else high.insert(x);
        rebalance();
    };
    auto erase = [&](int x) {
        auto it = low.find(x);
        if (it != low.end()) low.erase(it);
        else high.erase(high.find(x));
        rebalance();
    };

    for (int i = 0; i < k; i++) insert(a[i]);
    cout << *low.rbegin() << " ";
    for (int i = k; i < n; i++) {
        insert(a[i]); erase(a[i - k]);
        cout << *low.rbegin() << " ";
    }
    cout << "\n";
    return 0;
}
\end{lstlisting}

\begin{complexitybox}
\textbf{Time Complexity.} $O(n\log k)$ -- each insert/erase/rebalance is
$O(\log k)$. \textbf{Space Complexity.} $O(k)$.
\end{complexitybox}

\begin{takeawaybox}
The two-balanced-multisets structure keeps the median at the split point
under both insertion and deletion -- the erasable-heap the sliding window
needs. It directly powers the next problem, sliding window cost, by adding
running sums to each half.
\end{takeawaybox}
